\chapter{Newtonian Cosmology and the Friedmann Equation}

We were talking last time about the metric of space. The next question is dynamics: how does the scale factor change with time? First comes a little geometry, and then Newtonian cosmology. The Friedmann equation will not yet be derived from general relativity. It will be derived from Newton's equations.

\section{Homogeneity, Isotropy, and a Time-Dependent Metric}

Homogeneity and isotropy are the standing assumptions. If the universe is homogeneous and isotropic and it changes with time, then Hubble's law is always true. At any fixed time the recession velocity is proportional to distance,
\begin{equation}
v = H(t)\,d.
\end{equation}
The proportionality is linear. What is constant is not time, but space: the same proportionality factor applies everywhere at a given instant.

\begin{classroomqa}
\audiencequestion{Does that mean linearly proportional?}\lecturetimestamp{00:02:06}
\lecturerresponse{Yes. At a fixed time the recession velocity is linearly proportional to distance. The proportionality factor is the Hubble constant. In general it depends on time, but it does not vary from point to point in space.}
\end{classroomqa}

That is why the phrase ``Hubble constant'' is dangerous. It is constant with respect to position in space. It is not constant with respect to time.

A natural question comes next. If space expands, why do meter sticks not expand too? The answer is that meter sticks are held together by forces enormously stronger than anything associated with the cosmological expansion on atomic scales. Bound objects stay bound. What expands is the distance between sufficiently disconnected things, such as galaxies far enough apart that they do not hold each other together strongly.

\begin{classroomqa}
\audiencequestion{Is there any issue about the fact that the measuring rods are expanding too?}\lecturetimestamp{00:04:00}
\lecturerresponse{No. The measuring rods are not expanding. The coordinates are expanding. The forces which hold meter sticks together are far stronger than the effects due to the expansion, so ordinary bound objects do not expand with the background. Galaxies far enough apart do separate, but meter sticks do not.}
\end{classroomqa}

One can even imagine a tiny extra repulsive effect between atoms. That still would not tear a rod apart. It would only shift the equilibrium separation a little bit, the way a small extra pull changes the equilibrium position of a stiff spring.

\begin{classroomqa}
\audiencequestion{You can't model it though as a force between points, is that correct? Because otherwise it'd be an acceleration.}\lecturetimestamp{00:05:29}
\lecturerresponse{A small extra force between atoms can be imagined, but if it is tiny compared with the binding forces it only shifts the equilibrium position slightly. Expansion by itself does not create an ordinary force law. Some kinds of expansion can be modeled in terms of a force, namely accelerated expansion, but that is a later question.}
\end{classroomqa}

So the useful coordinates are coordinates embedded in the expanding matter distribution. Galaxies with no extra peculiar motion keep the same coordinate values. What changes with time is the relation between coordinate separation and physical distance measured in meters. That relation is codified by the scale factor \(a(t)\).

For flat space the spatial metric may be written
\begin{equation}
ds^2 = a(t)^2\left(dx^2+dy^2+dz^2\right).
\end{equation}
Here \(a\) is the actual distance, measured in meters, between neighboring points separated by one unit of coordinate distance. If the universe expands, \(a\) becomes time-dependent.

To pass to spacetime, start from the special-relativistic form
\begin{equation}
d\tau^2 = dt^2 - ds^2,
\end{equation}
and substitute the spatial metric. This gives
\begin{equation}
d\tau^2 = dt^2 - a(t)^2\left(dx^2+dy^2+dz^2\right).
\end{equation}
The coefficients do not depend on position in space. They depend only on time. That is the statement of homogeneity in metric language. The diagonal metric has the form
\begin{equation}
g_{\mu\nu}=\mathrm{diag}\!\left(1,-a(t)^2,-a(t)^2,-a(t)^2\right).
\end{equation}
\editorialnote{The full diagonal matrix written here follows the spoken metric-tensor description around 00:14:33. The board frame shows the line element clearly, but not a fully legible matrix.}

There are no off-diagonal terms. No \(dx\,dy\), no \(dx\,dt\), no other mixed terms. The metric is almost as simple as that of flat spacetime. The peculiarity is that the relation between spatial measure and time is itself time-dependent.

\begin{figure}
\centering
\includegraphics[width=0.62\textwidth]{lecture_02_figure_02.png}
\caption{Comoving grid sketch and spacetime interval with a time-dependent scale factor.}
\end{figure}

The picture carries the point. Alice and Bob remain one unit apart in the comoving coordinates, but that one coordinate unit represents more and more meters as time goes on. If \(a(t)\) were zero at some time in the past, then every physical separation would be zero. Everything would be squashed on top of everything else. That is why \(a\to 0\) corresponds to infinite density. If \(a\) grows, the universe becomes more and more diffuse.

This form of the metric is not valid in every frame. It is the form in the frame moving with the galaxies.

\begin{classroomqa}
\audiencequestion{In that equation, the small \(t\), \(a(t)\), whose \(t\) is that?}\lecturetimestamp{00:19:19}
\lecturerresponse{It is the time measured by clocks moving with the galaxies. Imagine identical clocks made in the same cosmic clock factory and carried by the galaxies. In this frame the time variable is exactly what those clocks measure. Along a galaxy worldline the proper time is the same as this coordinate time.}
\end{classroomqa}

Indeed, along a comoving galaxy worldline,
\begin{equation}
dx=dy=dz=0,
\end{equation}
so
\begin{equation}
d\tau^2 = dt^2,
\end{equation}
and therefore
\begin{equation}
dt=d\tau.
\end{equation}

\begin{classroomqa}
\audiencequestion{Are there transformations that leave this form of the metric invariant?}\lecturetimestamp{00:23:08}
\lecturerresponse{Rotations do, because \(dx^2+dy^2+dz^2\) is invariant under rotations. But Lorentz boosts do not preserve this form. Matter itself picks out the frame moving with the galaxies, and in some other frame the world would not look isotropic.}
\end{classroomqa}

\section{A Projectile in a Gravitational Field}

Before applying Newton to the whole universe, it helps to step back to a simpler problem: a heavy mass \(M\), a much smaller mass \(m\), and an initial outward impulse. After the impulse, the small mass simply moves in the gravitational field and does whatever Newton says it should do.

For the cosmological discussion, the important thing is not the full orbit problem. It is energy conservation. The kinetic energy is
\begin{equation}
\frac{1}{2}mv^2.
\end{equation}

\begin{figure}
\centering
\includegraphics[width=0.42\textwidth]{lecture_02_figure_03.png}
\caption{Outgoing test mass and the kinetic-energy term \(\frac{1}{2}mv^2\).}
\end{figure}

The gravitational potential energy contains the product of the masses, Newton's constant, and one power of the distance in the denominator. In the board writing the factors appear in the order \(mMG/D\), so it is best to keep that order here:
\begin{equation}
-\frac{mMG}{D}.
\end{equation}
Thus the total energy is
\begin{equation}
\frac{1}{2}mv^2-\frac{mMG}{D}.
\end{equation}
Energy conservation says only that this quantity does not change with time for the particular particle under discussion. Different particles may have different energies.

\begin{figure}
\centering
\includegraphics[width=0.62\textwidth]{lecture_02_figure_04.png}
\caption{Total Newtonian energy for a radially moving test mass: kinetic term minus gravitational potential.}
\end{figure}

The negative sign matters. Gravitational potential energy is negative because gravity is attractive. Bringing the particle closer creates a larger amount of negative potential energy. Moving it farther away makes the potential energy less negative.

\begin{classroomqa}
\audiencequestion{The potential energy increase with distance rather than decrease?}\lecturetimestamp{00:34:13}
\lecturerresponse{Yes. Forget the constants and look at the graphs of \(1/D\) and \(-1/D\). Force pushes toward decreasing potential energy. An attractive force therefore corresponds to the upside-down curve, \(-1/D\), so the potential energy increases with distance and becomes more negative as the particle moves inward.}
\end{classroomqa}

\begin{figure}
\centering
\includegraphics[width=0.46\textwidth]{lecture_02_figure_05.png}
\caption{Qualitative comparison of \(1/D\) and \(-1/D\) for the sign of gravitational potential energy.}
\end{figure}

\section{Escape Velocity}

Now classify the possibilities. The energy can be positive, negative, or zero. Those are the only three cases.

\begin{itemize}
\item If the energy is negative, the particle is bound. Throw it up with too small a velocity and it rises, slows, turns around, and falls back.
\item If the energy is positive, the particle escapes. Gravity decelerates it at first, but not enough to stop it, and eventually it moves off with a nonzero asymptotic outward velocity.
\item If the energy is exactly zero, the particle is at the escape velocity. That is the marginal case between escaping and not escaping.
\end{itemize}

For sufficiently small initial velocity, the potential term dominates from the start and the energy is negative.

\begin{figure}
\centering
\includegraphics[width=0.62\textwidth]{lecture_02_figure_06.png}
\caption{At small outward speed, the negative potential term dominates the total energy.}
\end{figure}

In the zero-energy case the distance keeps increasing, but the speed keeps decreasing. The particle cannot come to rest at any finite time. If it did, it would fall back in. So the distance is unbounded, while the velocity tends toward zero.

This is the direct analogy that matters for cosmology. The universe can behave the same way. If the initial outward Hubble flow is large enough, everything escapes from everything else. If it is too small, gravity wins and there is a giant crunch. There is also the marginal case, exactly analogous to escape velocity.

\section{Newton's Theorem}

The cosmological reduction works because of Newton's theorem for the inverse-square law.

If matter is distributed spherically symmetrically about some center, then the gravitational force on a particle outside the distribution is exactly the same as if all the mass were concentrated at the center. But the theorem is stronger than that. If the particle is inside the distribution, then only the mass at smaller radius contributes. The mass at larger radius gives no net force.

The cleanest version of the theorem is the thin shell. A perfect spherical shell exerts no gravitational force on a particle anywhere inside it, even away from the center. That is surprising at first sight, but it is true for the \(1/r^2\) law.

\begin{classroomqa}
\audiencequestion{In the example you gave there, should we be a little surprised that it seems like as the particle moves across the boundary, the force is discontinuous?}\lecturetimestamp{00:48:19}
\lecturerresponse{Yes. For an ideal infinitely thin shell it is surprising, and it is special to the \(1/r^2\) force law. A real shell has thickness, so the variation is smoothed out, but the sharp statement for an ideal shell is exactly the remarkable property of Newton's theorem.}
\end{classroomqa}

A picture then makes the theorem plausible. Replace a point mass by a point source of water. In three dimensions the outward flow decreases like \(1/r^2\). Put such sources continuously all around a spherical shell. Outside the shell the flow looks exactly as if everything were concentrated at the center. Inside the shell there can be no net flow, because there is no source at the center. The same structure is what underlies the gravitational theorem.

The point here is not to rederive the theorem. It is to use it.

\section{A Homogeneous Universe as a One-Body Problem}

Now go back to cosmology. Take a little patch of the world, filled with a uniform gas of galaxies. One may put oneself at the center if one likes. That is only a convenience. Anybody can call themselves at the center, and the same equation must come out.

Choose a galaxy on the \(x\)-axis at comoving coordinate \(x\). Since the universe is isotropic, the choice of axis does not matter. Draw a sphere centered on \(x=0\) and passing through that galaxy. By Newton's theorem, the force on the galaxy is exactly the same as if all the mass inside that sphere were concentrated at the center, while everything outside contributed nothing.

The sphere expands with the galaxies. It always contains the same galaxies, and therefore the same total mass, even though its physical radius changes with time.

The motion of that galaxy therefore obeys exactly the same kind of energy equation as the projectile:
\begin{equation}
\frac{1}{2}mv^2-\frac{mMG}{D}=\text{constant}.
\end{equation}
Here the \(M\) now means the mass inside the expanding sphere, not the central mass in the warm-up picture. The algebra is then simplified by multiplying by \(2\) and dividing by \(m\). Since a constant divided by \(m\) is still just a constant, one may write
\begin{equation}
v^2-\frac{2MG}{D}=\text{constant}.
\end{equation}
Again, ``constant'' means constant in time for that particular galaxy.

Now relate this to the scale factor. The physical distance to the galaxy is
\begin{equation}
D=ax,
\end{equation}
because \(x\) is a fixed comoving coordinate. Its velocity is
\begin{equation}
v=\dot D=\dot a\,x.
\end{equation}

The mass inside the sphere is the density times the volume. If \(\rho\) is the mass density, then
\begin{equation}
M=\frac{4\pi}{3}\rho D^3=\frac{4\pi}{3}\rho a^3x^3.
\end{equation}
Substituting into the energy equation gives
\begin{equation}
\dot a^2x^2-\frac{2G}{ax}\left(\frac{4\pi}{3}\rho a^3x^3\right)=\text{constant},
\end{equation}
or
\begin{equation}
\dot a^2x^2-\frac{8\pi G}{3}\rho a^2x^2=\text{constant}.
\end{equation}

Both terms on the left-hand side are proportional to \(x^2\). Everything else there is the same everywhere. Therefore the right-hand side must also be proportional to \(x^2\). That is not a contradiction with calling it a constant. It means constant with respect to time for each fixed \(x\), but it may still scale with \(x^2\) from one galaxy to another. That makes sense. The more distant galaxies are moving faster.

So divide away the common factor of \(x^2\). The remaining constant is now independent of both time and space. Write it in the historically standard way as \(-k\):
\begin{equation}
\dot a^2-\frac{8\pi G}{3}\rho a^2=-k.
\end{equation}
After dividing by \(a^2\),
\begin{equation}
\left(\frac{\dot a}{a}\right)^2=\frac{8\pi G}{3}\rho-\frac{k}{a^2}.
\end{equation}

This is the famous Friedmann-Robertson-Walker equation, though the lecture also remarks that it is really Friedman's equation. The sign convention is historical. Writing \(-k\) does not mean \(k\) must be positive. \(k\) may be positive, negative, or zero. In units with the speed of light set equal to \(1\), \(k\) is dimensionless.

\section{Density and the Matter-Dominated Equation}

The next question is immediate: what is \(\rho\)? It is not a constant. As the universe expands, the density decreases.

To express that dependence, take a unit cube in coordinate space: one unit on a side in \(x\), \(y\), and \(z\). The same galaxies remain inside that cube for all time because the cube expands with the coordinates. So the mass inside that coordinate cube is fixed.

\editorialnote{The lecture reuses the symbol \(M\) at several stages. Here \(M_{\Box}\) is used only to keep the mass in a unit coordinate cube distinct from the earlier central mass and sphere-enclosed mass.}

The physical side length of the cube is \(a\), so its physical volume is \(a^3\). Therefore
\begin{equation}
\rho=\frac{M_{\Box}}{a^3}.
\end{equation}
This is the matter-dominated case: the right-hand side involves only the density of ordinary, nonrelativistic matter.

\begin{classroomqa}
\audiencequestion{According to that equation, the Hubble constant is a function of \(\rho\). Therefore it's not a constant.}\lecturetimestamp{01:19:21}
\lecturerresponse{Exactly. The Hubble constant is not a constant. It depends on the density, and the density changes with time. The lecture even jokes that the final examination could consist of one question: Is the Hubble constant a constant?}
\end{classroomqa}

Substituting \(\rho=M_{\Box}/a^3\) into Friedmann's equation gives
\begin{equation}
\left(\frac{\dot a}{a}\right)^2=\frac{8\pi G M_{\Box}}{3a^3}-\frac{k}{a^2}.
\end{equation}

At this point the arbitrary choices are clear. One needs the mass in a unit coordinate cube and the value of \(k\). The three possible signs of \(k\) correspond to the three energy cases of the projectile problem.

\begin{itemize}
\item If \(-k>0\), the energy is effectively above escape velocity, and the universe expands forever.
\item If \(-k<0\), gravity eventually wins and the universe recollapses.
\item If \(k=0\), the universe is exactly at the escape velocity.
\end{itemize}

\begin{classroomqa}
\audiencequestion{Is \(k=0\)? Then the velocity between two galaxies is approaching zero?}\lecturetimestamp{01:21:12}
\lecturerresponse{Yes. In the critical case the velocity approaches zero, but the distance is still unbounded. The galaxies keep separating, only more and more slowly. If they actually came to rest at any finite time, they would fall back in.}
\end{classroomqa}

\section{The Critical Case \(k=0\)}

Now choose the simplest case, \(k=0\). That is the exact escape-velocity case, and it is simpler only because the mathematics is simpler. The equation becomes
\begin{equation}
\left(\frac{\dot a}{a}\right)^2=\frac{8\pi G M_{\Box}}{3a^3}.
\end{equation}

The right-hand side is always positive and never zero. Therefore \(\dot a/a\) cannot pass through zero. The universe cannot turn around in this case. It is either always expanding or always shrinking. Take the expanding branch.

Now guess a power law,
\begin{equation}
a(t)=ct^p.
\end{equation}
Then
\begin{equation}
\dot a = cp\,t^{p-1},
\end{equation}
so
\begin{equation}
\frac{\dot a}{a}=\frac{p}{t}.
\end{equation}
Substituting into the differential equation gives
\begin{equation}
\frac{p^2}{t^2}=\frac{8\pi G M_{\Box}}{3c^3t^{3p}}.
\end{equation}
For the powers of \(t\) to match, one must have
\begin{equation}
3p=2.
\end{equation}
So
\begin{equation}
p=\frac{2}{3}.
\end{equation}
Then the constants must satisfy
\begin{equation}
\frac{4}{9}=\frac{8\pi G M_{\Box}}{3c^3},
\end{equation}
hence
\begin{equation}
c^3=6\pi G M_{\Box}.
\end{equation}
Therefore
\begin{equation}
a(t)=\left(6\pi G M_{\Box}\right)^{1/3}t^{2/3}.
\end{equation}

There can still be an additive time shift. It is absorbed by replacing \(t\) with \(t-t_0\). Once one chooses the origin of time so that \(a=0\) at \(t=0\), that remaining constant is gone.

So the universe expands as the two-thirds power of time,
\begin{equation}
a(t)\propto t^{2/3}.
\end{equation}
It does not expand linearly with time. Linear growth would be the analog of unaccelerated motion. Here the expansion slows down.

If one pushes the formula backward to \(t=0\), the scale factor goes to zero and the density blows up. Saying \(a=0\) is extreme. We do not know what happens to matter when it becomes too dense. The real point is that at some early stage everything was so crowded that atoms, nuclei, and quarks were all packed together in an extraordinarily dense state.

\section{The Hubble Constant and the Age of the Universe}

In the critical model,
\begin{equation}
H(t)=\frac{\dot a}{a}=\frac{2}{3t}.
\end{equation}
So the Hubble constant is approximately one over the age of the universe,
\begin{equation}
t=\frac{2}{3H}.
\end{equation}
This is one of the main payoffs of the calculation. Measure the Hubble constant today, and one gets an age estimate for the universe, at least back to the point where the Newtonian model ceases to make sense.

The quoted age is roughly \(13\) to \(14\) billion years. The text also says that \(k=0\) is observationally very realistic: to the level of precision quoted there, no significant deviation from \(k=0\) is known.

\begin{classroomqa}
\audiencequestion{I'm just sort of confused about what \(M\) and \(a\) are.}\lecturetimestamp{01:39:16}
\lecturerresponse{Both contain arbitrariness from the choice of grid. Change the fiducial box and the amount of matter in it changes, and the numerical value of the scale factor changes too. What is physical is the density, \(M/a^3\), or equivalently combinations such as \(a/M^{1/3}\) that do not depend on how the grid is chosen. In the \(k=0\) solution the unknown constant multiplying \(t^{2/3}\) cancels out of \(\dot a/a\), so the age relation \(t=2/3H\) does not depend on the fiducial box mass.}
\end{classroomqa}

That cancellation is immediate in the solution. If
\begin{equation}
a(t)=Ct^{2/3},
\end{equation}
then
\begin{equation}
\dot a=\frac{2}{3}Ct^{-1/3},
\end{equation}
and therefore
\begin{equation}
\frac{\dot a}{a}=\frac{2}{3t}.
\end{equation}
The constant \(C\) drops out completely.

\section{Why the Newtonian Derivation Works}

Near the end the question becomes: why should the Newtonian derivation agree with the relativistic one as well as it does?

\begin{classroomqa}
\audiencequestion{This all seems to depend on an assumption of an open universe.}\lecturetimestamp{01:45:44}
\lecturerresponse{At this stage the derivation is Newtonian, and Newtonian physics means flat space. The interpretation of \(k\) in terms of the geometry of space is postponed to the next lecture. There \(k=0\) will correspond to flat space, while the other signs will be tied to curved spatial geometry.}
\end{classroomqa}

\begin{classroomqa}
\audiencequestion{Sometime long ago, did the length of meter sticks depend on \(a\)?}\lecturetimestamp{01:47:53}
\lecturerresponse{Far enough back, matter would have been so crowded that separate meter sticks would not really exist. Before that they would be squeezed together. The lecture also reminds us that relativity does not allow truly rigid bodies anyway: if a body were perfectly rigid, a push on one side would send a signal to the other side instantaneously, faster than light.}
\end{classroomqa}

\begin{classroomqa}
\audiencequestion{Why do you think there is this coincidence of equations, Newtonian and relativistic?}\lecturetimestamp{01:49:03}
\lecturerresponse{Because the Newtonian treatment is being applied only to a sufficiently small region. As long as the region is small compared with the radius of curvature of the whole universe, and as long as local Hubble velocities there are well below the speed of light, Newtonian mechanics is a good approximation. The lecture compares this with studying the growth of grass on the Earth: one can understand the local process without using the curvature of the whole Earth.}
\end{classroomqa}

This is the logic of the local patch. Newtonian mechanics assumes, first, that the curvature of space is not important for the question at hand, and second, that nothing locally is moving close to the speed of light. In cosmology that means choosing a region large enough to be homogeneous, but still small enough that curvature effects are negligible and Hubble velocities remain nonrelativistic.

With that in mind, the Newtonian calculation is not an accident. It is the local approximation to an equation that can also be read in relativistic form. If the curvature term is moved to the left, the Friedmann equation becomes
\begin{equation}
\left(\frac{\dot a}{a}\right)^2+\frac{k}{a^2}=\frac{8\pi G}{3}\rho.
\end{equation}
The left-hand side is geometric: it involves the scale factor and, next time, the curvature of space. The right-hand side is the energy density. That is exactly the kind of equation one expects from general relativity.

Next time the subject is curvature itself, Einstein's equations, and the relation between that relativistic derivation and the Newtonian one.
